\begin{objectivebox}[이 장에서 붙잡을 질문]
같은 네 비트 덧셈의 답을 유지하면서, 더 많은 논리를 써서 가장 늦게 도착하는 올림을 앞당길 수 있을까?
\end{objectivebox}

\section{올림 하나가 네 자리를 지난다}

3장에서는 네 전가산기를 연결해 \(0101_2+0011_2=1000_2\)을 계산했다. 그때는 각 자리의 합과 올림이 무엇인지에 집중했고, 결과가 언제 준비되는지는 세지 않았다. 이제 입력을 바꾸어 올림이 지나가는 길을 더 선명하게 보자.

\[
\begin{array}{r@{\ }l}
  &1111_2 \quad (15)\\[-1mm]
{}+&0001_2 \quad (1)\\
\hline
1\,&0000_2 \quad (16)
\end{array}
\]

가장 낮은 0번 자리에서는 \(1+1+0=10_2\)이므로 합 \(S_0=0\), 올림 \(c_1=1\)이 된다. 그 위의 세 자리에는 모두 \(A_i=1\), \(B_i=0\)이 놓여 있다. 따라서 아래에서 \(1\)이 들어올 때마다 \(1+0+1=10_2\)이 되어 현재 자리에는 \(0\)이 남고 올림 \(1\)은 다시 왼쪽으로 넘어간다. 마지막에는 네 합 비트 \(0000\)과 최종 올림 \(c_4=1\)을 얻는다.

이 과정을 매번 세 비트의 덧셈으로만 적는 대신, 한 자리가 올림에 어떤 일을 하는지 두 질문으로 나눌 수 있다.

\begin{description}[style=nextline,leftmargin=2.5em,font=\sffamily\bfseries]
  \item[생성 \(g_i\)] 이 자리의 두 입력만으로 새 올림을 만드는가?
  \item[전파 \(p_i\)] 아래 자리에서 올림이 들어오면 그것을 위 자리로 보내는가?
\end{description}

\keyterm{생성}(generate)은 두 입력이 모두 \(1\)인 경우다. 따라서 \(g_i=A_i\mathbin{\land}B_i\)로 둔다. \keyterm{전파}(propagate)는 두 입력이 서로 달라, 들어온 올림이 나가는 올림을 결정하는 경우다. 따라서 \(p_i=A_i\mathbin{\mathrm{XOR}}B_i\)로 둔다. 그러면 자리 \(i\)의 나가는 올림은

\[
  c_{i+1}=g_i\mathbin{\lor}(p_i\mathbin{\land}c_i)
\]

가 된다. 자체적으로 올림을 생성했거나, 두 입력이 서로 다른 자리에서 아래의 올림을 전파하면 다음 올림이 \(1\)이라는 뜻이다. 합 비트는

\[
  S_i=p_i\mathbin{\mathrm{XOR}}c_i
\]

로 얻는다.

\(1111_2+0001_2\)에서는 0번 자리가 \(g_0=1\)로 올림을 생성한다. 1·2·3번 자리는 \(p_1=p_2=p_3=1\)로 그 올림을 차례로 전파한다. 네 자리를 표로 확인하면 다음과 같다.

\begin{center}
\small
\begin{tabular}{cccccccc}
\toprule
\sffamily\bfseries 자리 \(i\) &
\sffamily\bfseries \(A_i\) &
\sffamily\bfseries \(B_i\) &
\sffamily\bfseries \(p_i\) &
\sffamily\bfseries \(g_i\) &
\sffamily\bfseries \(c_i\) &
\sffamily\bfseries \(S_i\) &
\sffamily\bfseries \(c_{i+1}\) \\
\midrule
0 & 1 & 1 & 0 & 1 & 0 & 0 & 1 \\
1 & 1 & 0 & 1 & 0 & 1 & 0 & 1 \\
2 & 1 & 0 & 1 & 0 & 1 & 0 & 1 \\
3 & 1 & 0 & 1 & 0 & 1 & 0 & 1 \\
\bottomrule
\end{tabular}
\end{center}

\begin{bookfigure}{올림 하나가 네 자리 안을 지나간다}{fig:ch04-carry-through-four-bits}
\begin{tikzpicture}[diagram]
  \path[use as bounding box] (0,0) rectangle (13.0,6.25);

  \node[diagram note,anchor=west,text=black!62] at (0.20,5.88)
    {높은 자리};
  \node[diagram note,anchor=east,text=black!62] at (12.80,5.88)
    {낮은 자리};
  \node[diagram note,anchor=center] at (6.50,5.88)
    {한 덧셈기 내부의 값 전달};

  % adder:carry-example
  \node[concept box,minimum width=26mm,minimum height=27mm] (fa3) at (1.70,3.95)
    {\textbf{자리 3}\\[0.8mm]\(A_3=1,\ B_3=0\)\\\(p_3=1,\ g_3=0\)\\[0.8mm]\(S_3=0\)};
  % adder:carry-example
  \node[concept box,minimum width=26mm,minimum height=27mm] (fa2) at (4.90,3.95)
    {\textbf{자리 2}\\[0.8mm]\(A_2=1,\ B_2=0\)\\\(p_2=1,\ g_2=0\)\\[0.8mm]\(S_2=0\)};
  % adder:carry-example
  \node[concept box,minimum width=26mm,minimum height=27mm] (fa1) at (8.10,3.95)
    {\textbf{자리 1}\\[0.8mm]\(A_1=1,\ B_1=0\)\\\(p_1=1,\ g_1=0\)\\[0.8mm]\(S_1=0\)};
  % adder:carry-example
  \node[concept emphasis,minimum width=26mm,minimum height=27mm] (fa0) at (11.30,3.95)
    {\textbf{자리 0}\\[0.8mm]\(A_0=1,\ B_0=1\)\\\(p_0=0,\ g_0=1\)\\[0.8mm]\(S_0=0\)};

  \draw[path wire] ([yshift=4mm]fa0.west)
    -- node[above,diagram note] {\(c_1=1\)} ([yshift=4mm]fa1.east);
  \draw[path wire] ([yshift=4mm]fa1.west)
    -- node[above,diagram note] {\(c_2=1\)} ([yshift=4mm]fa2.east);
  \draw[path wire] ([yshift=4mm]fa2.west)
    -- node[above,diagram note] {\(c_3=1\)} ([yshift=4mm]fa3.east);
  \draw[path wire] ([yshift=4mm]fa3.west)
    -- ++(-0.62,0) node[above,diagram note] {\(c_4=1\)};
  \draw[path wire] ([xshift=5.5mm,yshift=4mm]fa0.east)
    -- node[above,diagram note] {\(c_0=0\)} ([yshift=4mm]fa0.east);

  \node[concept emphasis,minimum width=94mm,minimum height=11mm] at (6.50,1.63)
    {\(1111_2+0001_2=1\,0000_2\) \quad \(g_0\)가 만들고 \(p_1,p_2,p_3\)가 전달한다};
  \node[diagram note,align=center,text=black!62] at (6.50,0.62)
    {이 입력은 값의 이동을 보여 준다 · 가장 긴 구조적 9층 경로는 그림 4.2에서 따로 센다};
\end{tikzpicture}
\end{bookfigure}


\figref{fig:ch04-carry-through-four-bits}은 이 입력에서 \(c_1,c_2,c_3,c_4\)가 모두 \(1\)이 되는 값의 연결을 보여 준다. 이것만으로 각 신호가 실제로 몇 초 뒤에 도착한다고 말할 수는 없다. 또한 회로의 가장 긴 구조적 경로는 어떤 입력값을 넣느냐에 따라 생겼다 사라지는 선이 아니다. 이 예는 그 긴 경로가 왜 필요한지 눈에 보이게 할 뿐이다.

최종 \(c_4=1\)은 이 무부호 계산의 정확한 답 \(16\)이 하위 네 비트의 범위 \(0\sim15\)를 벗어났음을 알려 준다. 그러나 3장에서 확인했듯 carry-out을 일반적인 부호 있는 넘침 신호라고 부르면 안 된다. 이 장의 비교 대상은 수의 부호가 아니라 \(c_4\)를 만드는 논리 경로다.

\section{비교하려면 무엇을 한 단위로 셀지 정해야 한다}

빠르다는 말을 숫자로 비교하려면 먼저 단위를 정해야 한다. 여기서는 실제 칩의 나노초나 전력을 추정하지 않는다. 다음과 같은 \keyterm{단위 게이트 지연}(unit gate delay) 모형을 사용한다.

\begin{itemize}
  \item XOR, AND, OR는 모두 입력이 두 개인 게이트로만 만들며, 게이트 하나를 수 1·지연 1층으로 센다.
  \item 배선 지연, 팬아웃, 입력 부하와 게이트 종류별 실제 지연 차이는 세지 않는다.
  \item 네 자리의 \(p_i\)와 \(g_i\)는 서로 독립이므로 1층에서 동시에 만든다.
  \item 여러 입력을 곱하는 AND는 2입력 게이트의 균형 트리로 만든다.
  \item 중간 곱항은 서로 공유하지 않는다.
\end{itemize}

비교 범위도 좁혀야 한다. 완성된 4비트 덧셈기의 네 합 비트 전체가 아니라, 가장 왼쪽의 최종 올림 \(c_4\)를 만드는 회로만 센다. 리플 방식에서 \(c_4\)는 네 자리의 정보를 모두 기다리는 출력이므로 가장 긴 올림 경로를 비교하기 좋다. 이 선택 때문에 뒤의 게이트 수는 완성된 덧셈기 전체의 크기가 아니다.

여기서 세는 층은 \emph{한 번의 덧셈 안에서 신호가 지나가는 게이트 경로}다. 8장에서 여러 덧셈 자체를 일렬 또는 나무 모양으로 배치하며 셀 의존 단계와는 층위가 다르다. 같은 ‘기다림’이라는 말을 쓰더라도 회로 안의 게이트 경로와 여러 계산 사이의 순서를 섞지 않는다.

\section{방법 A: 올림을 한 자리씩 전달한다}

가장 직접적인 방법은

\[
  c_{i+1}=g_i\mathbin{\lor}(p_i\mathbin{\land}c_i)
\]

를 0번 자리부터 3번 자리까지 차례로 적용하는 것이다. 아래 자리의 올림이 물결처럼 다음 자리로 이어진다고 해서 \keyterm{리플 캐리}(ripple carry)라고 부른다\parencite{mit6896fastaddition,mit6111arithmetic}.

\(p_i,g_i\)를 만드는 데 XOR 네 개와 AND 네 개, 모두 8게이트가 필요하다. 각 자리의 carry cell에는 \(p_i\land c_i\)를 만드는 AND 하나와, 그것을 \(g_i\)와 합치는 OR 하나가 들어간다. 네 자리에는 \(2\times4=8\)게이트가 필요하다. 따라서 이 모형에서 \(c_4\) 회로는

\[
  8+8=16\text{게이트}
\]

다.

가장 긴 구조적 경로는 \(p_i,g_i\)를 만드는 1층에서 시작한다. 그 뒤 네 carry cell이 각각 AND와 OR의 2층을 더한다. 따라서

\[
  1+4\times2=9\text{층}
\]

이다. 비트 폭을 \(n\)으로 바꾸면 같은 방식의 가장 긴 경로는 \(2n+1\)층으로 늘어난다.

\begin{center}
\small
\begin{tabular}{cc}
\toprule
\sffamily\bfseries 비트 폭 & \sffamily\bfseries 가장 긴 단위 게이트 경로 \\
\midrule
4비트 & \(2\times4+1=9\)층 \\
8비트 & \(2\times8+1=17\)층 \\
16비트 & \(2\times16+1=33\)층 \\
\bottomrule
\end{tabular}
\end{center}

이 표는 이 책이 정한 단순 모형의 구조적 증가를 보여 준다. 실제 덧셈기의 시간이나 특정 제품의 성능을 나타내지 않는다. 장점은 같은 carry cell을 규칙적으로 반복하며 \(c_1,c_2,c_3,c_4\)를 자연스럽게 모두 얻는다는 점이다. 대가는 비트 폭이 커질수록 마지막 올림까지 이어지는 사슬도 길어진다는 점이다.

\section{방법 B: 마지막 올림의 조건을 미리 펼친다}

한 자리씩 기다리지 않고 \(c_4\)가 \(1\)이 되는 조건을 처음부터 펼칠 수도 있다. 먼저

\[
  c_1=g_0\lor p_0c_0
\]

이고, 이것을 다음 자리에 대입하면

\[
  c_2=g_1\lor p_1g_0\lor p_1p_0c_0
\]

가 된다. 같은 대입을 3번과 4번 자리까지 이어 가면

\[
\begin{aligned}
  c_4={}&g_3
  \lor p_3g_2
  \lor p_3p_2g_1\\
  &\lor p_3p_2p_1g_0
  \lor p_3p_2p_1p_0c_0
\end{aligned}
\]

를 얻는다. 여기서 기호를 나란히 쓴 \(p_3p_2\)는 AND를 뜻한다. 최종 올림은 3번 자리에서 직접 생성되거나, 더 낮은 자리에서 생성된 올림이 그 위의 모든 자리를 전파하거나, 처음 들어온 \(c_0\)가 네 자리를 모두 전파할 때 \(1\)이 된다. 이런 생각을 \keyterm{올림 미리보기}(carry lookahead)라고 부른다\parencite{mit6111arithmetic,mit6896fastaddition}.

다섯 곱항을 \(T_0,\ldots,T_4\)라고 부르자. 2입력 AND를 균형 트리로 만들면 각 항이 준비되는 층은 다음과 같다.

\begin{center}
\small
\begin{tabular}{ccc}
\toprule
\sffamily\bfseries 이름 & \sffamily\bfseries 곱항 & \sffamily\bfseries 준비되는 층 \\
\midrule
\(T_0\) & \(g_3\) & 1층 \\
\(T_1\) & \(p_3g_2\) & 2층 \\
\(T_2\) & \(p_3p_2g_1\) & 3층 \\
\(T_3\) & \(p_3p_2p_1g_0\) & 3층 \\
\(T_4\) & \(p_3p_2p_1p_0c_0\) & 4층 \\
\bottomrule
\end{tabular}
\end{center}

다섯 항은 같은 시각에 준비되지 않는다. 따라서 OR 네 개의 개수만 보고 깊이를 정할 수 없다. 이 책에서는 도착 시각을 고려해 다음 연결을 사용한다.

\[
\begin{aligned}
  U_0&=T_0\lor T_1,\\
  U_1&=T_2\lor T_3,\\
  U_2&=U_0\lor U_1,\\
  c_4&=U_2\lor T_4.
\end{aligned}
\]

\(U_0\)는 3층, \(U_1\)은 4층, \(U_2\)는 5층, \(c_4\)는 6층에 준비된다. 같은 OR 네 개를 쓰더라도 늦게 준비되는 \(T_4\)를 더 불리한 잎에 놓으면 7층이 될 수 있다. 회로 깊이는 게이트 수만이 아니라 어떤 출력이 어떤 입력을 기다리는지까지 보아야 한다.

게이트 수도 직접 세어 보자. \(p_i,g_i\)에는 앞과 같은 8게이트가 든다. \(T_0\)에는 추가 AND가 없고, \(T_1\)부터 \(T_4\)까지는 각각 \(1,2,3,4\)개의 2입력 AND가 필요하다.

\[
  0+1+2+3+4=10\text{게이트}.
\]

다섯 항을 하나로 합치는 데에는 2입력 OR가 4개 필요하다. 따라서 합계는

\[
  8+10+4=22\text{게이트}
\]

이고, 가장 긴 경로는 6층이다.

\begin{bookfigure}{같은 \(c_4\)를 만드는 사슬과 펼친 식}{fig:ch04-chain-vs-expanded}
\begin{tikzpicture}[diagram]
  \path[use as bounding box] (0,0) rectangle (13.0,11.75);

  % ------------------------------------------------------------------
  % A. Ripple carry: 8 p/g gates + 8 carry gates = 16, depth 9.
  % The p/g gates are grouped because they are all evaluated in layer 1.
  % ------------------------------------------------------------------
  \draw[diagram panel] (0.05,8.05) rectangle (12.95,11.65);
  \node[diagram panel title] at (0.28,11.35)
    {A. 리플 캐리 · 16게이트 · 최장 9층};
  \node[diagram note,anchor=east] at (12.70,11.35)
    {1층은 8개 묶음 · 이후 상자 하나 = 게이트 하나};

  \foreach \x/\label in {
    1.05/1층,2.55/2층,3.90/3층,5.25/4층,6.60/5층,
    7.95/6층,9.30/7층,10.65/8층,12.00/9층}
  {
    \node[depth label] at (\x,10.78) {\label};
    \draw[depth guide] (\x,8.42) -- (\x,10.68);
  }

  % gate:ripple-pg
  % gate:ripple-pg
  % gate:ripple-pg
  % gate:ripple-pg
  % gate:ripple-pg
  % gate:ripple-pg
  % gate:ripple-pg
  % gate:ripple-pg
  \node[gate,minimum width=14mm,minimum height=13mm] (r-pg) at (1.05,9.47)
    {\(p_i\): XOR \(\times4\)\\
     \(g_i\): AND \(\times4\)};

  % gate:ripple-carry
  \node[path gate,minimum width=10.5mm,minimum height=8mm] (r-a0) at (2.55,9.47)
    {\(a_0\)\\AND};
  % gate:ripple-carry
  \node[path gate,minimum width=10.5mm,minimum height=8mm] (r-c1) at (3.90,9.47)
    {\(c_1\)\\OR};
  % gate:ripple-carry
  \node[path gate,minimum width=10.5mm,minimum height=8mm] (r-a1) at (5.25,9.47)
    {\(a_1\)\\AND};
  % gate:ripple-carry
  \node[path gate,minimum width=10.5mm,minimum height=8mm] (r-c2) at (6.60,9.47)
    {\(c_2\)\\OR};
  % gate:ripple-carry
  \node[path gate,minimum width=10.5mm,minimum height=8mm] (r-a2) at (7.95,9.47)
    {\(a_2\)\\AND};
  % gate:ripple-carry
  \node[path gate,minimum width=10.5mm,minimum height=8mm] (r-c3) at (9.30,9.47)
    {\(c_3\)\\OR};
  % gate:ripple-carry
  \node[path gate,minimum width=10.5mm,minimum height=8mm] (r-a3) at (10.65,9.47)
    {\(a_3\)\\AND};
  % gate:ripple-carry
  \node[path gate,minimum width=10.5mm,minimum height=8mm] (r-c4) at (12.00,9.47)
    {\(c_4\)\\OR};

  \draw[path wire] (r-pg.east) -- (r-a0.west);
  \draw[path wire] (r-a0.east) -- (r-c1.west);
  \draw[path wire] (r-c1.east) -- (r-a1.west);
  \draw[path wire] (r-a1.east) -- (r-c2.west);
  \draw[path wire] (r-c2.east) -- (r-a2.west);
  \draw[path wire] (r-a2.east) -- (r-c3.west);
  \draw[path wire] (r-c3.east) -- (r-a3.west);
  \draw[path wire] (r-a3.east) -- (r-c4.west);

  \node[diagram footer] at (6.50,8.35)
    {\(a_i=p_i\land c_i\), \(c_{i+1}=g_i\lor a_i\): 직접 \(p_i,g_i\) 선은 생략하고 올림 경로만 그렸다};

  % ------------------------------------------------------------------
  % B. Expanded c4: 8 p/g + 10 AND + 4 OR = 22, depth 6.
  % Direct inputs such as g1 and c0 are written inside their destination
  % gates. Only dependencies between intermediate gates are drawn.
  % ------------------------------------------------------------------
  \draw[diagram panel] (0.05,0.10) rectangle (12.95,7.80);
  \node[diagram panel title] at (0.28,7.50)
    {B. 펼친 \(c_4\) 식 · 22게이트 · 최장 6층};
  \node[diagram note,anchor=east] at (12.70,7.50)
    {8 \(p\)·\(g\) + 10 AND + 4 OR};

  \foreach \x/\label in {1.15/1층,3.15/2층,5.25/3층,7.35/4층,9.45/5층,11.55/6층}
  {
    \node[depth label] at (\x,6.90) {\label};
    \draw[depth guide] (\x,0.82) -- (\x,6.80);
  }

  % gate:expanded-pg
  % gate:expanded-pg
  % gate:expanded-pg
  % gate:expanded-pg
  % gate:expanded-pg
  % gate:expanded-pg
  % gate:expanded-pg
  % gate:expanded-pg
  \node[gate,minimum width=18mm,minimum height=43mm] (e-pg) at (1.15,4.10)
    {\(p_0,\ldots,p_3\)\\XOR \(\times4\)\\[1.2mm]
     \(g_0,\ldots,g_3\)\\AND \(\times4\)\\[1.2mm]
     \(T_0=g_3\)};

  % Six layer-2 AND gates.
  % gate:expanded-and
  \node[gate,minimum width=16mm] (e-t1) at (3.15,6.08) {\(T_1=p_3g_2\)};
  % gate:expanded-and
  \node[path gate,minimum width=16mm] (e-t2a) at (3.15,4.85) {\(p_3p_2\)};
  % gate:expanded-and
  \node[gate,minimum width=16mm] (e-t3a) at (3.15,4.28) {\(p_3p_2\)};
  % gate:expanded-and
  \node[gate,minimum width=16mm] (e-t3b) at (3.15,3.56) {\(p_1g_0\)};
  % gate:expanded-and
  \node[gate,minimum width=16mm] (e-t4a) at (3.15,2.58) {\(p_3p_2\)};
  % gate:expanded-and
  \node[gate,minimum width=16mm] (e-t4b) at (3.15,1.86) {\(p_1p_0\)};

  % Layer 3: one OR and three AND gates.
  % gate:expanded-or
  \node[gate,minimum width=17mm,minimum height=7mm] (e-u0) at (5.25,6.08)
    {\(U_0\)\\\(T_0\lor T_1\)};
  % gate:expanded-and
  \node[path gate,minimum width=18mm,minimum height=7mm] (e-t2) at (5.25,4.85)
    {\(T_2\)\\\((p_3p_2)g_1\)};
  % gate:expanded-and
  \node[gate,minimum width=19mm,minimum height=7mm] (e-t3) at (5.25,3.92)
    {\(T_3\)\\\((p_3p_2)(p_1g_0)\)};
  % gate:expanded-and
  \node[gate,minimum width=19mm,minimum height=7mm] (e-t4c) at (5.25,2.22)
    {\((p_3p_2)\)\\\(\cdot(p_1p_0)\)};

  % Layer 4: U1 and T4 become ready independently.
  % gate:expanded-or
  \node[path gate,minimum width=17mm,minimum height=7mm] (e-u1) at (7.35,4.85)
    {\(U_1\)\\\(T_2\lor T_3\)};
  % gate:expanded-and
  \node[gate,minimum width=17mm,minimum height=7mm] (e-t4) at (7.35,2.22)
    {\(T_4\)\\\(\cdots\land c_0\)};

  % Layer 5 and 6: finish the arrival-time-aware OR tree.
  % gate:expanded-or
  \node[path gate,minimum width=17mm,minimum height=7mm] (e-u2) at (9.45,4.85)
    {\(U_2\)\\\(U_0\lor U_1\)};
  % gate:expanded-or
  \node[path gate,minimum width=17mm,minimum height=7mm] (e-c4) at (11.55,4.85)
    {\(c_4\)\\\(U_2\lor T_4\)};

  % Inputs from layer 1 to the six layer-2 gates. These are short,
  % separated horizontal segments, so no wire crosses a gate or label.
  \draw[wire] ([yshift=19.8mm]e-pg.east) -- (e-t1.west);
  \draw[path wire] ([yshift=7.5mm]e-pg.east) -- (e-t2a.west);
  \draw[wire] ([yshift=1.8mm]e-pg.east) -- (e-t3a.west);
  \draw[wire] ([yshift=-5.4mm]e-pg.east) -- (e-t3b.west);
  \draw[wire] ([yshift=-15.2mm]e-pg.east) -- (e-t4a.west);
  \draw[wire] ([yshift=-22.4mm]e-pg.east) -- (e-t4b.west);

  % T0 arrives directly from layer 1 and is written inside U0. Drawing its
  % long bypass wire would obscure the layer-to-layer flow.
  \draw[wire] (e-t1.east) -- (e-u0.west);
  \draw[path wire] (e-t2a.east) -- (e-t2.west);
  \draw[wire] (e-t3a.east) -- ([yshift=1.4mm]e-t3.west);
  \draw[wire] (e-t3b.east) -- ([yshift=-1.4mm]e-t3.west);
  \draw[wire] (e-t4a.east) -- ([yshift=1.4mm]e-t4c.west);
  \draw[wire] (e-t4b.east) -- ([yshift=-1.4mm]e-t4c.west);

  \draw[path wire] (e-t2.east) -- (e-u1.west);
  \draw[wire] (e-t3.east) -- ++(0.12,0) |- ([yshift=-1.4mm]e-u1.west);
  \draw[wire] (e-t4c.east) -- (e-t4.west);
  \draw[wire] (e-u0.east) -- ++(2.25,0) |- ([yshift=1.4mm]e-u2.west);
  \draw[path wire] (e-u1.east) -- (e-u2.west);
  \draw[path wire] (e-u2.east) -- (e-c4.west);
  \draw[wire] (e-t4.east) -- ++(2.30,0) |- ([yshift=-1.4mm]e-c4.west);

  \node[diagram footer] at (6.50,0.50)
    {굵은 선·회색 상자 = 최장 6층 경로 · 1층에서 직접 오는 입력은 도착 게이트 안에 적었다};
\end{tikzpicture}
\end{bookfigure}


\figref{fig:ch04-chain-vs-expanded}의 위쪽은 16게이트와 9층의 리플 사슬을, 아래쪽은 22게이트와 6층의 펼친 식을 보여 준다. 굵은 선은 각 회로에서 가장 긴 경로다. 아래 회로에서 \(p_3p_2\) 같은 중간 곱항을 여러 항이 공유하지 않고 다시 만든 것은 게이트 수를 재현 가능하게 고정하기 위한 선택이다.

\section{더 짧은 길은 더 많은 논리와 배선을 요구한다}

두 회로는 같은 입력에서 같은 \(c_4\)를 만들어야 한다. 달라지는 것은 정답이 아니라 그 정답에 이르는 회로량과 가장 긴 경로다.

\begin{center}
\small
\begin{tabular}{lcc}
\toprule
\sffamily\bfseries \(c_4\)를 만드는 교육 회로 &
\sffamily\bfseries 게이트 수 &
\sffamily\bfseries 가장 긴 단위 게이트 경로 \\
\midrule
리플 캐리 & 16 & 9층 \\
펼친 미리보기 식 & 22 & 6층 \\
\bottomrule
\end{tabular}
\end{center}

펼친 식은 게이트를 6개 더 쓰는 대신 가장 긴 경로를 3층 줄였다. 그러나 이 비교는 의도적으로 비대칭이다. 리플 회로는 계산 도중 \(c_1,c_2,c_3\)도 자연스럽게 내지만, 22게이트 회로는 \(c_4\) 하나만 만들도록 세었다. \(c_1\)부터 \(c_4\)까지 각각 펼치고 곱항을 전혀 공유하지 않으면

\[
  8+(1+3+6+10)+(1+2+3+4)=38\text{게이트}
\]

가 필요하다. 첫 8은 모든 \(p_i,g_i\), 가운데 합은 각 carry의 AND, 마지막 합은 OR의 수다. 네 합 비트 \(S_i\)를 만드는 XOR까지 넣으면 다시 수가 달라진다.

실제 빠른 덧셈기는 이 두 교육 회로 가운데 하나만 그대로 고르지 않는다. 중간 항을 공유하거나, 비트를 묶어 계층을 만들거나, 접두 구조·carry-select·carry-bypass 같은 다른 방법을 사용할 수 있다\parencite{mit6004designtradeoffs,mit6111arithmetic}. 어떤 구조도 모든 비트 폭과 배선 조건에서 유일한 정답은 아니다. 여기서 배워야 할 일반 원리는 하나다. 같은 답을 유지해도 더 많은 논리와 배선을 제공하면 가장 긴 의존 경로를 줄일 수 있다.

그렇다고 게이트 수와 경로 길이를 하나의 점수로 더할 수는 없다. 게이트 수는 회로량을 가늠하는 교육 단위이고, 경로의 층수는 이 모형에서 기다림을 가늠하는 단위다. 실제 면적·전력·시간에는 게이트 종류, 배선, 팬아웃, 공정이 함께 작용한다. 이 장의 숫자는 선택의 방향을 드러내기 위한 비교값이지 제품 사양이 아니다.

\section{다음에는 더 많은 부분 결과를 합쳐야 한다}

3장에서는 한 자리의 합과 올림을 만들었고, 이 장에서는 여러 자리의 올림을 어떤 순서로 계산할지 비교했다. 이로써 같은 덧셈도 회로를 어떻게 펼치고 묶느냐에 따라 기다리는 길이가 달라진다는 사실을 보았다.

다음 장의 곱셈에서는 한 줄의 올림보다 더 많은 중간값이 생긴다. 각 자리에서 부분곱을 만들고, 여러 부분곱을 다시 더해야 한다. 그 결과를 한 줄씩 차례로 합칠 수도 있고, 더 많은 회로를 써서 여러 줄을 함께 줄일 수도 있다. 덧셈에서 만난 질문이 더 큰 모습으로 돌아오는 셈이다.

\section*{이번 장의 선택과 비용}
\addcontentsline{toc}{section}{이번 장의 선택과 비용}

\begin{center}
\small
\begin{tabular}{>{\raggedright\arraybackslash\sffamily\bfseries}p{0.29\textwidth} >{\raggedright\arraybackslash}p{0.58\textwidth}}
\toprule
보존한 기능·정확한 결과 & \(1111_2+0001_2=1\,0000_2\), 최종 올림 \(c_4=1\) \\
\midrule
계산이 요구하는 양 & 같은 \(p_i,g_i\) 값과 같은 carry 식, 최종 올림이 의존하는 네 자리 정보 \\
\midrule
기계가 제공할 자원 & \(c_4\) 기준 리플 16게이트·9층, 펼친 미리보기 22게이트·6층 \\
\midrule
유한 정밀도에서 달라질 수 있는 것 & 두 회로 모두 하위 네 비트 \(0000\)과 \(c_4=1\)을 냄; 무부호 범위는 넘지만 구조 선택이 signed overflow의 의미를 바꾸지는 않음 \\
\midrule
설계 대가 & 추가 논리와 배선을 제공해 가장 긴 단위 게이트 경로를 줄임 \\
\midrule
아직 해결되지 않은 제약 & 덧셈보다 많은 부분 결과를 합치는 곱셈에서는 어떤 순서가 필요한가 \\
\bottomrule
\end{tabular}
\end{center}
